% Figura 11.3: Analogía estructural de opciones reales: VAN estático vs Opción Call de diferimiento (Caso Telefónica UGR)
\begin{tikzpicture}[
    font=\small\sffamily,
    col/.style={
        rectangle, rounded corners=4pt, draw=#1, fill=#1!8,
        line width=1.1pt, minimum width=5.2cm, minimum height=0.82cm,
        align=center, font=\footnotesize\sffamily, text=dfinslate
    }
]

% Cabeceras de correspondencia financiera vs real
\node[rectangle, rounded corners=3pt, fill=dfinnavy, text=white, minimum width=5.2cm, minimum height=0.82cm,
      font=\footnotesize\sffamily\bfseries] (h1) at (-3.4, 5.2) {Opción Call Financiera ($C_0$)};
\node[rectangle, rounded corners=3pt, fill=dfinemerald, text=white, minimum width=5.2cm, minimum height=0.82cm,
      font=\footnotesize\sffamily\bfseries] (h2) at (3.4, 5.2) {Opción Real de Inversión (Caso Telefónica)};

% 5 Filas pareadas de Black-Scholes / Opciones Reales
\node[col=dfinnavy] (f1) at (-3.4, 4.0) {\textbf{$S_0$:} Precio actual de la acción};
\node[col=dfinemerald] (r1) at (3.4, 4.0) {\textbf{$\text{VA}_0$:} Valor actual de flujos del proyecto};

\node[col=dfinnavy] (f2) at (-3.4, 2.95) {\textbf{$E$:} Precio de ejercicio pactado};
\node[col=dfinemerald] (r2) at (3.4, 2.95) {\textbf{$I_0$:} Desembolso necesario para desplegarlo};

\node[col=dfinnavy] (f3) at (-3.4, 1.9) {\textbf{$T$:} Tiempo hasta el vencimiento};
\node[col=dfinemerald] (r3) at (3.4, 1.9) {\textbf{$T$:} Horizonte antes de perder la licencia/ventaja};

\node[col=dfinnavy] (f4) at (-3.4, 0.85) {\textbf{$\sigma$:} Volatilidad del subyacente};
\node[col=dfinemerald] (r4) at (3.4, 0.85) {\textbf{$\sigma$:} Incertidumbre de la demanda o tecnología};

\node[col=dfinnavy] (f5) at (-3.4, -0.2) {\textbf{$r_f$:} Tipo de interés libre de riesgo};
\node[col=dfinemerald] (r5) at (3.4, -0.2) {\textbf{$r_f$:} Coste del dinero sin riesgo en el plazo $T$};

% Flechas dobles de isomorfismo
\foreach \i in {1,2,3,4,5} {
    \draw[<->, >=Latex, line width=1pt, color=dfincopper] (f\i.east) -- (r\i.west);
}

% Mini gráfica inferior comparando VAN rígido (S - E, recta negativa si S < E) vs Opción Real (max(S-E,0) + valor temporal)
\begin{axis}[
    at={(-5.4cm, -5.2cm)},
    width=11.4cm, height=5.0cm,
    xlabel={Valor presente de los flujos operativos ($S = \text{VA}_{\text{flujos}}$)},
    ylabel={Valor neto hoy},
    xmin=40, xmax=160, ymin=-45, ymax=105,
    axis x line=bottom, axis y line=left,
    xtick={60, 100, 140},
    xticklabels={{60}, {$I_0 = E = 100$}, {140}},
    ytick={-30, 0, 30, 60, 90},
    grid=major, grid style={dotted, color=dfinmuted!40},
    tick label style={font=\scriptsize\sffamily, color=dfinslate},
    label style={font=\scriptsize\sffamily\bfseries, color=dfinnavy},
    legend style={at={(0.03,0.97)}, anchor=north west, font=\scriptsize\sffamily,
                  fill=white, draw=dfinslate!35, rounded corners=2pt}
]
\addplot[forget plot, domain=40:160, samples=2, color=dfinslate!55, line width=0.7pt] {0};
% VAN estático rígido (decisión irreversible hoy sin esperar): S - 100
\addplot[domain=55:155, samples=2, color=dfinburg, line width=1.3pt, dashed] {x - 100};
\addlegendentry{VAN rígido hoy ($S - I_0$, sin flexibilidad)}

% Payoff intrínseco en T: max(S - 100, 0)
\addplot[domain=45:155, samples=80, color=dfinnavy, line width=1.45pt] {max(x - 100, 0)};
\addlegendentry{Valor intrínseco $\max(S - E,\ 0)$}

% Valor vivo de la opción real con flexibilidad de esperar
\addplot[domain=45:155, samples=80, color=dfinemerald, line width=1.65pt]
    { 0.5*((x-100) + sqrt((x-100)^2 + 324)) };
\addlegendentry{VAN estratégico = VAN básico $+$ Valor de la Opción Real}

\node[font=\scriptsize\sffamily\bfseries, text=dfinemerald, anchor=south east, fill=white, inner sep=1.5pt, align=right]
    at (axis cs:97, 8) {Valor temporal de esperar\\[-1pt]($>0$ aunque $S < I_0$)};

\end{axis}

\end{tikzpicture}
